%-------------------------------------------------------------------------------------------------
% LaTeX Resume for Professionals
% Author : Rodrigo Zanini
% License : CC BY 4.0 https://creativecommons.org/licenses/by/4.0/
%-------------------------------------------------------------------------------------------------
\documentclass{article}

\usepackage[T1]{fontenc}
\usepackage[sfdefault,light]{FiraSans}
\usepackage{fontawesome}
\usepackage{xcolor}
\usepackage{calc}
\usepackage{graphicx}
\usepackage{tikz}
\graphicspath{ {./} }

\usepackage[cm,headings]{fullpage}

\usepackage[inline]{enumitem}
\usepackage{tabularx}
\usepackage{fancyhdr}
\pagestyle{fancy}

\usepackage{hyperref}
\hypersetup{
	colorlinks=true,
	urlcolor=primary,
}
\usepackage{multirow}


% page numbering is disabled because whole document fits on a single page
\pagenumbering{gobble}

%-------------------------------------------------------------------------------------------------
%---------- CUSTOM COMMANDS
%-------------------------------------------------------------------------------------------------
%---------- Icons
\newcommand{\icon}[1]{\begin{tabular}{p{\iconWidth}}{#1}\end{tabular}}
\renewcommand{\section}[2]{\vspace{0.5em}\color{primary}\textbf{{#2}{#1}}\hrule\color{black}}

%---------- Head
\renewcommand{\headrulewidth}{0pt}
\newcommand{\cvText}[1]{#1\vspace{0.75em}}

%---------- Entries
% Entry with Name, Time, Position, Location and Details
\newcommand{\cvEntryNTPLD}[5]{
	\begin{itemize}[leftmargin=\entryMargin]
		\item[]
		\begin{tabularx}{\textwidth-\entryMargin}{Xr}
			\textbf{\color{black}#1} & {\color{secondary}\small#2} \\
			\textit{\color{secondary}\small#3} & \textit{\color{secondary}\small#4} \\
		\end{tabularx}\vspace{-0.75em}
		\begin{itemize}[labelsep=\entryMargin-1.5em,leftmargin=\entryMargin]
			\small#5
		\end{itemize}
	\end{itemize}
}
% Entry with Name, Time, Position and Location
\newcommand{\cvEntryNTPL}[4]{
	\begin{itemize}[leftmargin=\entryMargin]
		\item[]
		\begin{tabularx}{\textwidth-\entryMargin}{Xr}
			\textbf{\color{black}#1} & {\color{secondary}\small#2} \\
			\textit{\color{secondary}\small#3} & \textit{\color{secondary}\small#4} \\
		\end{tabularx}\vspace{-0.75em}
	\end{itemize}
}
% Entry with Name, Time, Position and Details
\newcommand{\cvEntryNTPD}[4]{
	\begin{itemize}[leftmargin=\entryMargin]
	\item[]
	\begin{tabularx}{\textwidth-\entryMargin}{Xr}
		\textbf{\color{black}#1} & {\color{secondary}\small#2} \\
		\textit{\color{secondary}\small#3} \\
	\end{tabularx}\vspace{-0.75em}
	\begin{itemize}[labelsep=\entryMargin-1.5em,leftmargin=\entryMargin]
		\small#4
	\end{itemize}
\end{itemize}
}
% Entry with Name, Time and Details
\newcommand{\cvEntryNTD}[3]{
	\begin{itemize}[leftmargin=\entryMargin]
		\item[]
		\begin{tabularx}{\textwidth-\entryMargin}{Xr}
			\textbf{\color{black}#1} & {\color{secondary}\small#2} \\
		\end{tabularx}\vspace{-0.75em}
		\begin{itemize}[labelsep=\entryMargin-1.5em,leftmargin=\entryMargin]
			\small#3
		\end{itemize}
	\end{itemize}
}
% Entry with Name and Details
\newcommand{\cvEntryND}[2]{
	\begin{itemize}[leftmargin=\entryMargin]
		\item[]
		\begin{tabularx}{\textwidth-\entryMargin}{Xr}
			\textbf{\color{black}#1} \\
		\end{tabularx}\vspace{-0.75em}
		\begin{itemize}[labelsep=\entryMargin-1.5em,leftmargin=\entryMargin]
			\small#2
		\end{itemize}
	\end{itemize}
}
% Entry with Name
\newcommand{\cvEntryN}[1]{
	\small
	\begin{itemize}[leftmargin=\entryMargin]
		#1
	\end{itemize}
	\normalsize
}

% Entry for Languages (up to 4)
\newcommand{\cvItemL}[8]{
	\item[]
	\begin{tabularx}{\textwidth-\entryMargin}{XXXX}
		{\textbf{#1} \textit{#2}}&{\textbf{#3} \textit{#4}}&{\textbf{#5} \textit{#6}}&{\textbf{#7} \textit{#8}}
	\end{tabularx}\vspace{-0.25em}
}
% Entry for Academical Projects
\newcommand{\cvEntryAcademicalNTD}[3]{\cvEntryNTD{\educationIcon{Учебный: }#1}{#2}{#3}}
\newcommand{\cvEntryAcademicalND}[2]{\cvEntryND{\educationIcon{Учебный: }#1}{#2}}

% Entry for Personal Projects
\newcommand{\cvEntryPersonalNTD}[3]{\cvEntryNTD{\personalIcon{Личный: }#1}{#2}{#3}}
\newcommand{\cvEntryPersonalND}[2]{\cvEntryND{\personalIcon{Личный: }#1}{#2}}

% Entry for Work Projects
\newcommand{\cvEntryWorkNTD}[3]{\cvEntryNTD{\workIcon{Рабочий: }#1}{#2}{#3}}
\newcommand{\cvEntryWorkND}[2]{\cvEntryND{\workIcon{Рабочий: }#1}{#2}}

% Bullet point under an entry's details list
\newcommand{\cvItem}[1]{\item[\bulletIcon]{#1\vspace{-0.25em}}}
% Special bullet point under an entry's details list
\newcommand{\cvItemS}[1]{\item[\accentIcon]{#1\vspace{-0.25em}}}
% Bulletless point under an entry's details list, with Name and Description 
\newcommand{\cvItemND}[2]{\item[]{\textbf{#1}\hspace{1em}#2}\vspace{-0.25em}}
% Simple bulletless text under an entry's details list
\newcommand{\cvItemT}[1]{\item[]{#1\vspace{-0.25em}}}

% List of skills under an entry's details list
\newcommand{\cvItemSkillset}[1]{\item[]{\begin{linelist}#1\end{linelist}\vspace{-0.25em}}}
% link to github project
\newcommand{\ghlink}[1]{\small\githubIcon\hspace*{-0.4em}\href{https://github.com/#1}{#1}}
\newcommand{\cvItemGithubLink}[1]{\cvItemT{\ghlink{#1}}}

% custom single-line list with items separated with bullet
\newcommand\sbullet[1][.5]{\mathbin{\vcenter{\hbox{\scalebox{#1}{$\bullet$}}}}}
\newlist{linelist}{enumerate*}{1}
\setlist[linelist]{label={}, afterlabel={}, itemjoin={{ $\sbullet$ }}}

%-------------------------------------------------------------------------------------------------
%---------- SETTINGS HERE
%-------------------------------------------------------------------------------------------------
%---------- Colours
\definecolor{primary}{HTML}{7d26e0}
\definecolor{secondary}{HTML}{3C454D}
\definecolor{info}{HTML}{17A2B8}

%---------- Margins
\raggedbottom
\raggedright
\setlength{\tabcolsep}{0in}
\setlength{\voffset}{-0.5cm}
\setlength{\headheight}{3.5em}
\addtolength{\headsep}{-2em}
\addtolength{\oddsidemargin}{-0.25cm}
\addtolength{\evensidemargin}{-0.25cm}
\addtolength{\headwidth}{0.5cm}
\addtolength{\textwidth}{0.5cm}

%---------- Entries
\def \entryMargin{1em}

%---------- Icons
\def \iconWidth{1.5em}
% Predefined icons, based on FontAwesome.
% see https://ctan.mirror.rafal.ca/fonts/fontawesome/doc/fontawesome.pdf for more options.
\def \linkedinIcon{\icon{\faLinkedin}}
\def \telegramIcon{\icon{\faSend}}
\def \phoneIcon{\icon{\faPhone}}
\def \homeIcon{\icon{\faHome}}
\def \emailIcon{\icon{\faEnvelope}}
\def \githubIcon{\icon{\faGithub}}
\def \websiteIcon{\icon{\faGlobe}}

\def \educationIcon{\icon{\faGraduationCap}}
\def \workIcon{\icon{\faBriefcase}}
\def \projectsIcon{\icon{\faFlask}}
\def \communicationIcon{\icon{\faComments}}
\def \awardsIcon{\icon{\faTrophy}}
\def \skillsIcon{\icon{\faGears}}
\def \interestsIcon{\icon{\faGamepad}}

\def \expertIcon{\icon{\faStar}}
\def \proficientIcon{\icon{\faStarHalfFull}}
\def \noviceIcon{\icon{\faStarO}}
\def \personalIcon{\icon{\faUser}}

\def \bulletIcon{\icon{\faAngleRight}}
\def \accentIcon{\icon{\faAngleDoubleRight}} % \faCaretRight \faAngleDoubleRight \faCode


%-------------------------------------------------------------------------------------------------
%---------- DATA HERE
%-------------------------------------------------------------------------------------------------
%---------- Header data
\def \name{Владимир Редькин}
\def \nameSuffix{}
\def \subtitleText{Python/Rust разработчик, SRE}
% 
\def \summaryText{summary if I had one}

\def \linkedinLink{https://www.linkedin.com/}
\def \linkedinText{/vsredkin}

\def \phoneText{+7 111 111-11-11}

\def \homeText{Казань, республика Татарстан}

\def \emailLink{mailto:vsredkin01@gmail.com}
\def \emailText{vsredkin01@gmail.com}

\def \githubLink{https://github.com/hectorhw}
\def \githubText{/HectorHW}

\def \websiteLink{website link}
\def \websiteText{website text}

\def \telegramLink{https://t.me/hectorhw}
\def \telegramText{@hectorhw}

%---------- Header format
\def \fullName{\textbf{\huge\name}\large\hspace{0.3em}\textit{\nameSuffix}}
\def \subtitle{\textit{\small\cvText{\subtitleText}}}
\def \summary{\cvText{\summaryText}}
\def \linkedin{\small\linkedinIcon\href{\linkedinLink}{\linkedinText}}
\def \phone{\small\phoneIcon{\phoneText}}
\def \home{\small\homeIcon{\homeText}}
\def \email{\small\emailIcon\href{\emailLink}{\emailText}}
\def \github{\small\githubIcon\hspace*{-0.3em}\href{\githubLink}{\githubText}}
\def \website{\small\websiteIcon\href{\websiteLink}{\websiteText}}
\def \telegram{\small\telegramIcon\href{\telegramLink}{\telegramText}}


%-------------------------------------------------------------------------------------------------
%---------- START
%-------------------------------------------------------------------------------------------------
\begin{document}
	%-------------------------------------------------------------------------------------------------
	%---------- HEADER
	%-------------------------------------------------------------------------------------------------
		\fancyhead[L]{
			\begin{tabular}[c]{ll}
				\multirow{2}{1.7cm}{
					\begin{tikzpicture}[overlay, remember picture]
						\clip (0.6,0.3)  circle (0.85cm) ;
						\node[anchor=center, %anchor is upper left corner of the graphic
						xshift=1.97cm, %shifting around
						yshift=-0.88cm] 
						at (current page.north west) %left upper corner of the page
						{\includegraphics[width=2.3cm]{photo}}; 
					\end{tikzpicture}
				} & {\fullName} \\
				& {\subtitle}
			\end{tabular}
			\vspace{0.5cm}
		}
		\fancyhead[R]{
			\begin{tabular}[c]{l@{\hspace{1em}}l@{\hspace{1em}}l}
				% Configure the order in which the header data appears. Must be in 3 colums.
				 {\telegram} & {\github} & {\email}\\
				 \multicolumn{3}{l}{\home}
				\vspace{0.5em}
			\end{tabular}
			\vspace{1em}
		}


	%-------------------------------------------------------------------------------------------------
	%---------- INTRODUCTION
	%-------------------------------------------------------------------------------------------------
	%\summary\\ % Comment out if not using
	%-------------------------------------------------------------------------------------------------
	%---------- EDUCATION
	%-------------------------------------------------------------------------------------------------
	\section{Образование}{\educationIcon}
	\cvEntryNTPD{Казанский федеральный университет}{2019 -- 2023}{Фундаментальная информатика (бакалавриат)}{
		\cvItemSkillset{
			\item Вышмат
			\item Теория алгоритмов
			\item Программная инженерия
			\item Базы данных
			\item Криптография
			\item Классический ML
		}
	}
	\cvEntryNTPD{Тинькофф Образование}{2021 -- 2023}{Академия Бэкенда}{
		\cvItemSkillset{
			\item Паттерны проектирования
			\item SOLID
			\item Проектирование API
			\item Сеть
			\item Базы данных
			\item Мониторинг
		}
	}
	\cvEntryNTPD{Казанский федеральный университет}{2019 -- 2021}{Переводчик в сфере профессиональной коммуникации}{
		\cvItemSkillset{
			\item Английский язык
			\item Письменный перевод
			\item Синхронный перевод
		}
	}
	
	%-------------------------------------------------------------------------------------------------
	%---------- WORK EXPERIENCE
	%-------------------------------------------------------------------------------------------------
	\section{Опыт работы}{\workIcon}
	\cvEntryNTPD
	{Тинькофф}{январь 2023 -- июль 2023}
	{Стажёр-разработчик (департамент ИТ-инфраструктуры)}{
		\cvItem{
			Чат-бот, отслеживающий новые релизы на GitHub
		}
		\cvItem{
			Визуализация инфраструктуры
		}
		\cvItem{
			
			Развитие системы генерации конфигурации prometheus
		}
	}
	
	%-------------------------------------------------------------------------------------------------
	%---------- PROJECTS
	%-------------------------------------------------------------------------------------------------
	%\newpage %Uncomment if FIRST ITEM of this section is not entirely on this page
	\section{Проекты}{\projectsIcon}
	\cvEntryWorkND{Чат-бот, отслеживающий новые релизы на GitHub}{
		\cvItemT{Бот выводит уведомления о новых версиях, прикрепляет changelog и при необходимости позволяет заводить задачки на обновление в Jira. Бот использует S3 для хранения состояния, развёрнут в k8s и конфигурируется через git.}
		\cvItemSkillset{
			\item Python
			\item Pydantic
			\item FastAPI
			\item S3
			\item Docker
			\item Kubernetes
			\item Helm
		}
	}

	\cvEntryWorkND{Визуализация инфраструктуры}{
		\cvItemT{
			Делал панель для мониторинга серверов. Панель содержит довольно большое число правил, для автоматической генерации написан скрипт, полагающийся на два jinja-темплейта.
		}
		\cvItemSkillset{
			\item Python
			\item Poetry
			\item Jinja
			\item Grafana
			\item FlowCharting
		}
	}

	\cvEntryWorkND{Развитие системы генерации конфигурации prometheus}{
		\cvItemT{
			На замену ранее существовавшему костыльному коду создан jinja-темплейт и разработан yaml-based DSL для описания устройств и автоматической генерации таргетов prometheus.
		}
		\cvItemSkillset{
			\item Python
			\item Poetry
			\item Jinja
			\item Pydantic
			\item PyYaml
		}
	}

	\cvEntryPersonalND{Homelab}{
		\cvItemGithubLink{HectorHW/nginx\_configurator}
		\cvItemT{
			Набор сервисов, развёрнутых на моём домашнем сервере + Cloudflare, для которых я потихоньку делаю автоматизацию и мониторинг.
		}
		\cvItemSkillset{
			\item Linux
			\item Nginx
			\item Cloudflare
			\item Prometheus
			\item GitHub Actions
		}
	}
	
	\cvEntryAcademicalND{Интерпретируемый язык программирования}{
		\cvItemGithubLink{HectorHW/blop-lang}
		\cvItemT{
			В рамках курсовой работы написал интерпретатор функционального языка программирования, включающий парсер, байткодовую виртуальную машину и сборщик мусора.
		}
		\cvItemSkillset{
			\item Rust
			\item PEG parsing
			\item Interpreter
			\item Language runtime
			\item GC
		}
	}

	\cvEntryAcademicalND{Грамматический фаззер}{
		\cvItemGithubLink{HectorHW/bocchifuzz}
		\cvItemT{
			В рамках дипломной работы разработал фаззер, позволяющий описывать входные данные в виде грамматики либо набора начальных образцов и проводить фазз-тестирование без модификации или перекомпиляции исследуемых файлов.
		}
		\cvItemSkillset{
			\item Rust
			\item PEG parsing
			\item TUI
			\item Ptrace
		}
	}

	%-------------------------------------------------------------------------------------------------
	%---------- INTERSESTS
	%-------------------------------------------------------------------------------------------------
	%\newpage %Uncomment if FIRST ITEM of this section is not entirely on this page
	\section{Остальные интересы и увлечения}{\interestsIcon}
	\cvEntryN{
		\cvItemND{Bio-inspired computing}{Люблю изучать вещи на стыке информатики и биологии: генетические алгоритмы, swarm intelligence, нейронные сети.}
		\cvItemND{Волонтёрство}{Помогаю с организацией научно-популярных лекций в Национальной библиотеке Республики Татарстан.}
	}
\end{document}